\subsection{Recursion with a fixed rational seed}
\label{subsec:seed-recursion}

Here we fix one rational seed, for instance that of
Corollary~\ref{cor:seed} or of Theorem~\ref{thm:bounded-seed}, and bound
two other costs of block recursion with it: the bit operations needed
for integer inputs (Proposition~\ref{prop:integer-bit-cost}) and a
precision that suffices for a prescribed forward error
(Proposition~\ref{prop:seed-error}). Let the seed have side
\(d_0\ge2\) and \(R\) products, and put \(p_{\rm seed}=\log_{d_0}R\);
this is not the pooled exponent \(p\). Write the seed as
\[
 \operatorname{vec}(XY)
 =\mathsf W\bigl((\mathsf U\operatorname{vec}X)\odot
   (\mathsf V\operatorname{vec}Y)\bigr),
 \qquad \mathsf U,\mathsf V\in\Q^{R\times d_0^2},\quad
 \mathsf W\in\Q^{d_0^2\times R},
\]
for \(d_0\times d_0\) matrices \(X\) and \(Y\), where \(\odot\) is the
entrywise product. \emph{Block recursion} multiplies \(n\times n\)
matrices by padding them with zeros to side \(d_0^\ell\), where
\(\ell=\lceil\log_{d_0}n\rceil\), applying this identity to the
\(d_0\times d_0\) arrays of blocks, computing the \(R\) block products
recursively, and cropping the result to its leading \(n\times n\) block.
Constants in \(O_{\rm seed}\) depend only on the seed.

\begin{lemma}[Block recursion with a rational seed]\label{lem:seed-recursion}
We have \(R>d_0^2\), and block recursion multiplies \(n\times n\)
matrices with at most \(c_{\rm seed}n^{p_{\rm seed}}\) arithmetic
operations, where \(c_{\rm seed}=(1+6Rd_0^2)R\).
\end{lemma}

\begin{proof}
The rows of \(\mathsf U\) are the left forms of the seed, linear forms
in the entries of \(X\). The flattening of \(\MM_{d_0}\) along leg~\(1\)
has rank \(d_0^2\) (Step~1 of the proof of Lemma~\ref{lem:batch}), and
its image lies in the span of the left forms, so \(R\ge d_0^2\). If
equality held, the left forms would be a basis of the dual space, so
some nonzero \(X\) would make \(\mathsf U\operatorname{vec}X\) the first
unit vector. Then \(\operatorname{vec}(XY)\) would be a multiple of the
first column of \(\mathsf W\) for every \(Y\), so the map \(Y\mapsto XY\)
would have rank at most one, whereas its rank is
\(d_0\operatorname{rank}(X)\ge d_0\ge2\). Hence \(R>d_0^2\).

At a node of side \(N\), a multiple of \(d_0\), forming the \(R\) left
and the \(R\) right block combinations costs at most
\(2R(2d_0^2-1)(N/d_0)^2\) operations, and forming the \(d_0^2\) output
blocks costs at most \(d_0^2(2R-1)(N/d_0)^2\); together this is less
than \(6RN^2\). Let \(\theta_{\ell'}\) be the cost at side
\(d_0^{\ell'}\). Then \(\theta_0=1\) and
\(\theta_{\ell'}\le R\theta_{\ell'-1}+6Rd_0^{2\ell'}\), so, using
\(R-d_0^2\ge1\),
\[
 \theta_\ell\le R^\ell+6Rd_0^{2\ell}\sum_{\ell'=0}^{\ell-1}
 \Bigl(\frac R{d_0^2}\Bigr)^{\ell'}
 \le(1+6Rd_0^2)R^\ell.
\]
Since \(d_0^\ell<d_0n\), we have
\(R^\ell=(d_0^\ell)^{p_{\rm seed}}<(d_0n)^{p_{\rm seed}}
=Rn^{p_{\rm seed}}\).
\end{proof}

For the integer bound, choose a positive integer \(c_{\rm den}\) that
clears all denominators of \(\mathsf U\), \(\mathsf V\) and \(\mathsf W\),
and put
\begin{gather*}
 c_{\rm coef}=\max\{1,\max|c_{\rm den}\mathsf U|,\max|c_{\rm den}\mathsf V|,
   \max|c_{\rm den}\mathsf W|\},\\
 c_{\rm in}=\lceil\log_2(d_0^2c_{\rm coef})\rceil,\qquad
 c_{\rm out}=\lceil\log_2(Rc_{\rm coef})\rceil,
\end{gather*}
where \(\max|\cdot|\) is the largest absolute entry.

\begin{proposition}[Integer bit complexity]\label{prop:integer-bit-cost}
Let \(n\ge2\), \(\ell=\lceil\log_{d_0}n\rceil\) and \(b_{\rm in}\ge1\). Block
recursion with the integer matrices \(c_{\rm den}\mathsf U\),
\(c_{\rm den}\mathsf V\), \(c_{\rm den}\mathsf W\), followed by one exact
division of every entry by \(c_{\rm den}^{3\ell}\), multiplies two
\(n\times n\) integer matrices whose entries
have absolute value below \(2^{b_{\rm in}}\) exactly, with intermediate
integers of bit length at most
\[
 2b_{\rm in}+\ell(2c_{\rm in}+c_{\rm out})+2,
\]
including a sign bit, and with
\(O_{\rm seed}\bigl(n^{p_{\rm seed}}(b_{\rm in}+\log n)^2\bigr)\) bit
operations, using schoolbook integer arithmetic. The cost of
constructing the seed is not included.
\end{proposition}

\begin{proof}
Use the integer matrices \(c_{\rm den}\mathsf U\), \(c_{\rm den}\mathsf V\),
\(c_{\rm den}\mathsf W\) at every level, without intermediate divisions.
One level then computes \(c_{\rm den}^3\) times the product, so by
induction on the depth the root computes \(c_{\rm den}^{3\ell}XY\);
divide every entry of the result by \(c_{\rm den}^{3\ell}\) at the end.
These divisions are exact.

The absolute row sums of \(c_{\rm den}\mathsf U\) and
\(c_{\rm den}\mathsf V\) are at most \(d_0^2c_{\rm coef}\), so each
input-transform level multiplies the largest absolute entry by at most
\(d_0^2c_{\rm coef}\), and a leaf product has absolute value below
\(2^{2b_{\rm in}}(d_0^2c_{\rm coef})^{2\ell}\). Each output-transform level
multiplies the bound by at most \(Rc_{\rm coef}\). Hence every
intermediate value, including partial sums and products by
coefficients, has absolute value below
\(2^{2b_{\rm in}}(d_0^2c_{\rm coef})^{2\ell}(Rc_{\rm coef})^\ell
\le2^{2b_{\rm in}+\ell(2c_{\rm in}+c_{\rm out})}\), so the stated bit length
suffices, with one bit to spare. Evaluating the one-level integer
identity on all-one inputs gives
\(c_{\rm den}^3d_0\le Rd_0^4c_{\rm coef}^3\), so the divisor
\(c_{\rm den}^{3\ell}\) also has this bit length. By
Lemma~\ref{lem:seed-recursion}, the recursion uses
\(O_{\rm seed}(n^{p_{\rm seed}})\) integer operations, and forming the
divisor and dividing take \(O(n^2+\ell)\) more. Since
\(\ell=O(\log n)\), each operation costs
\(O_{\rm seed}((b_{\rm in}+\log n)^2)\) bit operations.
\end{proof}

For entries with a bounded number of bits, the cost is
\(O_{\rm seed}(n^{p_{\rm seed}}\log^2n)\). For the seed of
Theorem~\ref{thm:bounded-seed} or Corollary~\ref{cor:seed} for a
rational \(\eps>0\), we have \(p_{\rm seed}<9/4+\eps\), so such matrices
are multiplied with \(O_\eps(n^{9/4+\eps})\) bit operations. Rational
inputs can be handled by first clearing one common denominator for each
matrix, with \(b_{\rm in}\) bounding the resulting integer entries; if the
\(n^2\) entries have unrelated \(b_{\rm in}\)-bit denominators, this common
denominator may need up to \(n^2b_{\rm in}\) bits.

For the precision bound, write \(\|Z\|_{\max}\) for the largest absolute
entry of a matrix \(Z\) and \(\|Z\|_\infty\) for its largest absolute
row sum, and put
\(\Upsilon=\|\mathsf U\|_\infty\|\mathsf V\|_\infty\|\mathsf W\|_\infty\).

\begin{proposition}[A sufficient precision]\label{prop:seed-error}
Let \(X\) and \(Y\) be \(n\times n\) matrices with \(n\ge2\), and put
\[
 \ell=\lceil\log_{d_0}n\rceil,\qquad
 \chi=\ell(2d_0^2+R+3)+1,\qquad
 \gamma_\chi=\frac{\chi u}{1-\chi u}.
\]
Compute \(XY\) by block recursion, evaluating every entry
\(\sum_kc_kw_k\) of each linear combination by recursive summation
\((\cdots((c_1w_1+c_2w_2)+c_3w_3)+\cdots)\) in a fixed order of \(k\),
and the leaves as scalar products. Suppose that each arithmetic operation and
each stored nonzero coefficient has relative error at most \(u\), that
zeros are stored exactly, and that no operation underflows or
overflows. If \(\chi u<1\), the computed product \(\widehat{XY}\)
satisfies
\begin{equation}\label{eq:seed-forward-error}
 \|\widehat{XY}-XY\|_{\max}
 \le\gamma_\chi\Upsilon^\ell\|X\|_{\max}\|Y\|_{\max}.
\end{equation}
Here \(X\) and \(Y\) are the stored inputs; errors in storing the
original inputs are not counted.
\end{proposition}

\begin{proof}
The idea is to expand each computed entry into bilinear monomials and to
count the rounding factors along one path of the recursion. Expand a
computed output entry into signed bilinear monomials in the entries of
the padded inputs, each carrying the factors \(1+\delta\) of the
operations that produced it. At each input-transform level, a monomial
picks up one factor from storing a coefficient, one from multiplying by
it, and at most \(d_0^2-1\) from additions, so its two input paths
contribute at most \(2\ell(d_0^2+1)\) factors. Each output-transform
level contributes at most \(R+1\) factors, and the leaf product one. The
total is at most \(\chi\), also when a stored coefficient is reused, and
a product of at most \(\chi\) factors \(1+\delta\) with \(|\delta|\le u\)
differs from one by at most \(\gamma_\chi\). Applying the three row-sum
bounds at each level shows that the absolute values of the exact
monomials sum to at most \(\Upsilon^\ell\|X\|_{\max}\|Y\|_{\max}\). The
triangle inequality proves \eqref{eq:seed-forward-error}.
\end{proof}

Applying the seed to all-one inputs shows \(\Upsilon\ge d_0\ge2\). Hence,
for \(0<\vartheta\le1\), the condition
\[
 u\le\frac{\vartheta}{2\chi\Upsilon^\ell}
\]
gives \(\chi u\le1/4\) and an error of at most
\(\vartheta\|X\|_{\max}\|Y\|_{\max}\). In binary floating point with
\(t_{\rm sig}\) significand bits, including the leading bit, the unit
roundoff is \(2^{-t_{\rm sig}}\) for round-to-nearest and at most
\(2^{1-t_{\rm sig}}\) for directed
rounding. Hence the condition holds with
\[
 \left\lceil\log_2\frac{2\chi}{\vartheta}\right\rceil
 +\ell\lceil\log_2\Upsilon\rceil+2
\]
significand bits, that is, with
\(O_{\rm seed}(\log n+\log(1/\vartheta))\) bits. \citet{DDHK07}, following \citet{BL80}, give a forward error
analysis of this kind for a large class of recursive matrix
multiplication algorithms; Proposition~\ref{prop:seed-error} states the
bound with explicit constants for one rational seed and dot products
evaluated from left to right.

\subsection{One algorithm for all sizes}
\label{subsec:uniform-setup}

The algorithms of Section~\ref{subsec:relative-schemes} depend on a
fixed allowance \(\eps\). Letting the allowance shrink with the input
size, and counting the preprocessing, gives one algorithm for all sizes.

\begin{corollary}[One algorithm for all sizes]\label{cor:uniform-setup}
There is one algorithm that, for every \(n\), multiplies two
\(n\times n\) matrices over a field of characteristic zero in
\(n^{9/4+o(1)}\) steps, where a step is one step of a fixed multitape
Turing machine during the preprocessing, or one addition, subtraction,
multiplication or division in the field on matrix entries or seed
coefficients.
\end{corollary}

\begin{proof}
The idea is to try the allowances \(1/v\), \(v=1,2,\ldots\), within a
preprocessing budget of \(n\) machine steps each, and to use the
smallest allowance that succeeds. At side \(n\), for each integer
\(1\le v\le\lfloor\log_2n\rfloor\), run a fixed Turing-machine
implementation of the procedure of
Theorem~\ref{thm:construction-resources} for the allowance \(1/v\) for
at most \(n\) steps; within the same budget compute \(c_{\rm seed}\) of
Lemma~\ref{lem:seed-recursion} for the resulting seed and test
\(c_{\rm seed}^v\le n\). Interrupt any attempt that exceeds its budget.
Use the seed of the largest successful \(v\), or classical
multiplication if there is none. Including binary counters for the
budgets, this preprocessing takes \(O(n(\log n)^2)\) steps. The
coefficients of the selected seed have total bit length at most \(n\),
so converting them to field elements by Horner's rule in base~\(2\) and
one division each takes \(O(n)\) field operations.

For every fixed \(v\), the procedure and the test finish within the
budget once \(n\) is large enough, and then \(c_{\rm seed}^v\le n\).
Hence the selected \(v=v(n)\) tends to infinity. Once some \(v\)
succeeds, Theorem~\ref{thm:bounded-seed} gives
\(p_{\rm seed}<9/4+1/v\), and Lemma~\ref{lem:seed-recursion} bounds the
multiplication by
\(c_{\rm seed}n^{p_{\rm seed}}\le n^{1/v}n^{9/4+1/v}=n^{9/4+2/v(n)}\)
operations, which proves the claim.
\end{proof}
